\section{Повторяющиеся игры}

В данном разделе рассматривается класс динамических игр, в которых одна и та же игра (называемая шаговой) повторяется несколько раз. В зависимости от горизонта планирования выделяют конечно и бесконечно повторяющиеся игры.

\subsection{Конечно повторяющиеся игры}

Если базовая игра имеет единственное равновесие, тривиальное повторение стратегий может казаться единственным исходом. Рассмотрим дилемму заключенного, которая повторяется дважды. Если применить обратную индукцию, то на втором шаге игроки будут играть равновесие Нэша шаговой игры (оба предадут). Следовательно, на первом шаге игра также сведется к однократной дилемме заключенного, и игроки предадут друг друга с самого начала.

\begin{definition}
Пусть $G$ — игра в нормальной форме, $T > 1$ — натуральное число. Обозначим за $G(T)$ динамическую игру, в которой игра $G$ повторяется $T$ раз, а выигрыш игрока $i$ составляет:
\begin{equation}
    u_i = \sum_{t=1}^T u_i(a_t),
\end{equation}
где $a_t \in A$ — действия игроков в момент времени $t = 1, \dots, T$. Игра $G$ называется шаговой игрой, а элементы множества $A$ — действиями.
\end{definition}

Анализ стратегий в повторяющихся играх требует учета информационных множеств. В дважды повторенной дилемме заключенного у каждого игрока 5 информационных множеств (одно на первом шаге и четыре на втором). Поскольку в каждом множестве доступно два хода, общее число стратегий составляет $2^5 = 32$.

Однако многие стратегии эквивалентны. Формально, стратегия $\sigma_i$ эквивалентна $\sigma_i'$, если для всех $j$ и профилей стратегий остальных игроков $\sigma_{-i}$ выполняется:
\begin{equation}
    u_j(\sigma_i, \sigma_{-i}) = u_j(\sigma_i', \sigma_{-i})
\end{equation}
С учетом эквивалентности у каждого игрока в данной игре остается 8 уникальных стратегий.

\begin{theorem}
Пусть $G = \langle I, A, u \rangle$ — игра в нормальной форме c единственным равновесием $a^*$, и $T > 1$. Тогда в игре $G(T)$ существует единственное совершенное по подыграм равновесие Нэша, такое, что в каждой подыгре, начинающейся с $t \ge 1$, игроки выберут действия $a^*$.
\end{theorem}

\begin{proof}
Доказательство проводится методом математической индукции.
\begin{itemize}
    \item Все подыгры $G(T)$, начинающиеся в последний момент времени $t = T$, полностью совпадают с шаговой игрой $G$. Следовательно, в них будет сыграно равновесие $a^*$.
    \item Все подыгры, начинающиеся с $t = T - 1$, совпадают с шаговой игрой $G$ с точностью до прибавления константы $u(a^*)$ к функциям выигрышей. Таким образом, равновесные действия в момент $t = T - 1$ также будут $a^*$.
    \item Согласно принципу индукции, данный шаг применим для всех $t = 1, \dots, T$.
\end{itemize}
\end{proof}

Ситуация меняется, если в шаговой игре существует несколько равновесий. Ожидания относительно поведения в будущем могут влиять на поведение в настоящем.

\begin{example}
Рассмотрим игру $G$, в которой матрица выигрышей имеет следующий вид:
\begin{table}[h]
\centering
\renewcommand{\arraystretch}{1.3}
\begin{NiceTabular}{cc|ccc}
\toprule
& & \multicolumn{3}{c}{\textbf{Игрок 2}} \\
& & A & B & C \\
\midrule
\multirow{3}{*}{\textbf{Игрок 1}} 
& A & $1, 1$ & $5, 0$ & $0, 0$ \\
& B & $0, 5$ & $4, 4$ & $0, 0$ \\
& C & $0, 0$ & $0, 0$ & $3, 3$ \\
\bottomrule
\end{NiceTabular}
\caption{Матрица выигрышей шаговой игры с множественными равновесиями}
\end{table}

В данной базовой игре два равновесия Нэша в чистых стратегиях: $(A, A)$ и $(C, C)$. Профиль $(B, B)$ равновесием не является.
При повторении игры дважды (в $t=1$ и $t=2$) возникает 9 подыгр на втором шаге. В момент времени $t=2$ игроки обязаны реализовывать равновесные стратегии шаговой игры: либо $(A, A)$, либо $(C, C)$. Причем профиль стратегий, сыгранных в $t=2$, может зависеть от исхода $t=1$. 

Предположим, игроки договариваются: в $t=2$ они реализуют $(C, C)$, только если в $t=1$ было реализовано $(B, B)$. Во всех остальных случаях в $t=2$ играется $(A, A)$. Тогда приведенная матрица выигрышей для всей игры с учетом этих стратегий примет вид:
\begin{table}[h]
\centering
\renewcommand{\arraystretch}{1.3}
\begin{NiceTabular}{cc|ccc}
\toprule
& & \multicolumn{3}{c}{\textbf{Игрок 2}} \\
& & A & B & C \\
\midrule
\multirow{3}{*}{\textbf{Игрок 1}} 
& A & $2, 2$ & $6, 1$ & $1, 1$ \\
& B & $1, 6$ & $7, 7$ & $1, 1$ \\
& C & $1, 1$ & $1, 1$ & $4, 4$ \\
\bottomrule
\end{NiceTabular}
\caption{Матрица выигрышей игры $G(2)$ при условных стратегиях}
\end{table}

В этой модифицированной игре появляется три равновесия: $(A, A)$, $(C, C)$ и $(B, B)$. Таким образом, профиль $(B, B)$, не являющийся равновесием в шаговой игре, становится частью равновесной стратегии в повторяющейся игре благодаря угрозе наказания на втором шаге.
\end{example}

\subsection{Бесконечно повторяющиеся игры}

Бесконечно повторяющиеся игры моделируют ситуации, в которых взаимодействие агентов продолжается неопределенно долго. Поскольку вероятность окончания игры после каждого шага строго меньше единицы (или агенты просто дисконтируют будущее), обратная индукция здесь неприменима. 

\begin{definition}
Пусть $G = \langle I, A, u \rangle$ — игра в нормальной форме, $\delta < 1$ — фактор дисконтирования (дисконт). Обозначим моменты времени за $t = 0, 1, \dots$, а действия игроков в момент $t$ за $a_t \in A$. Обозначим за $G(\delta)$ игру $G$, повторяемую бесконечное число раз, в которой выигрыш игрока $i$ равен усредненному дисконтированному платежу:
\begin{equation}
    \bar{u}_i(\delta) = (1 - \delta) \sum_{t=0}^\infty \delta^t u_i(a_t)
\end{equation}
\end{definition}

Нормирующий множитель $(1 - \delta)$ используется для того, чтобы масштаб выигрышей в бесконечной игре был сопоставим с выигрышами в однократной шаговой игре.

Стратегия в бесконечно повторяющейся игре — это план действий, зависящий от всей доступной предыстории. Поскольку предыстория может быть бесконечно долгой, вводится класс марковских стратегий.

\begin{definition}
Пусть $G = \langle I, A, u \rangle$ — игра в нормальной форме. Марковская стратегия с памятью $T$ для игрока $i$ есть функция:
\begin{equation}
    m_i : A^T \to A_i,
\end{equation}
которая определяет действие $a_{it}$ для каждой истории игры длиной $T$.
\end{definition}

Известными примерами таких стратегий в дилемме заключенного являются:
\begin{itemize}
    \item \textit{Всегда предавать} или \textit{Никогда не предавать} (память $T=0$).
    \item \textit{«Око за око»} (Tit-for-tat): повторять действие партнера из предыдущего периода.
    \item \textit{Триггерная стратегия} (Grim trigger): не предавать, пока партнер не предал ни разу. После первого же предательства со стороны партнера — всегда предавать.
\end{itemize}

\subsection{Народная теорема}

Народная теорема (Folk Theorem) формализует интуицию о том, что при достаточно терпеливых игроках (высоком дисконте $\delta$) в бесконечно повторяющейся игре можно поддержать практически любой исход в качестве равновесия.

\begin{example}
Рассмотрим дилемму заключенного:
\begin{table}[h]
\centering
\renewcommand{\arraystretch}{1.3}
\begin{NiceTabular}{cc|cc}
\toprule
& & \multicolumn{2}{c}{\textbf{Вася}} \\
& & \textit{«cooperate»} & \textit{«defect»} \\
\midrule
\multirow{2}{*}{\textbf{Петя}} 
& \textit{«cooperate»} & $3, 3$ & $0, 5$ \\
& \textit{«defect»} & $5, 0$ & $1, 1$ \\
\bottomrule
\end{NiceTabular}
\end{table}

Если оба игрока используют триггерные стратегии, усредненный выигрыш каждого равен $3$. Если в момент времени $t=0$ Игрок 1 решает отклониться и сыграть \textit{«defect»}, его выигрыш составит:
\begin{equation}
    \bar{u}_1' = (1 - \delta) \left( u_1(D, C) + \sum_{t=1}^\infty \delta^t u_1(D, D) \right) = (1 - \delta) \left( 5 + \frac{\delta}{1 - \delta} \cdot 1 \right)
\end{equation}
Отклонение невыгодно, если $\bar{u}_1' \le 3$, то есть:
\begin{equation}
    (1 - \delta) \left( 5 + \frac{\delta}{1 - \delta} \right) \le 3
\end{equation}
Это неравенство выполняется при $\delta \ge \frac{1}{2}$. Следовательно, при достаточно большом дисконте кооперация является равновесием.
\end{example}

Для формализации Народной теоремы введем математический аппарат. Пусть $G = \langle I, A, u \rangle$ — игра в стратегической форме, а $u(a)$ — вектор выигрышей всех игроков. Обозначим множество всех возможных выпуклых комбинаций выигрышей (кооперативное множество):
\begin{equation}
    X = \left\{ x \in \mathbb{R}^N \, \middle| \, x = \sum_{a \in A} \alpha_a u(a), \sum_{a \in A} \alpha_a = 1, \alpha_a \ge 0 \right\}
\end{equation}

Также определим минимаксный выигрыш — уровень полезности, который игрок $i$ может себе гарантировать, даже если все остальные игроки объединятся против него:
\begin{equation}
    \underline{u}_i = \min_{a_{-i} \in A_{-i}} \max_{a_i \in A_i} u_i(a_i, a_{-i})
\end{equation}
Выигрыш игрока $i$ в каждом периоде принципиально не может быть меньше $\underline{u}_i$. Обозначим множество выигрышей, строго превышающих минимаксные:
\begin{equation}
    Y = \{ y \in \mathbb{R}^N \mid y_i > \underline{u}_i \}
\end{equation}

\begin{theorem}[Народная теорема Нэша]
Пусть $G = \langle I, A, u \rangle$ — игра в стратегической форме. Пусть целевой вектор выигрышей $u \in X \cap Y$. Тогда существует такое $\bar{\delta} < 1$, что для всех $\delta \in [\bar{\delta}, 1)$ в повторяющейся игре $G(\delta)$ существует равновесие в чистых стратегиях, в котором приведенный выигрыш каждого игрока $i$ равен $u_i$.
\end{theorem}

\begin{proof}[Схема доказательства]
Определим профиль наказывающих действий $\mu_{-i}^i$, обеспечивающих игроку $i$ выигрыш не более минимаксного:
\begin{equation}
    \mu_{-i}^i = \arg\min_{a_{-i} \in A_{-i}} \max_{a_i \in A_i} u_i(a_i, a_{-i})
\end{equation}
Пусть $m = (m_1, \dots, m_N)$ — профиль марковских стратегий, дающий целевой выигрыш $u \in X \cap Y$. Построим стратегию: пока все играют $m$, продолжаем играть $m$. Если игрок $i$ единожды отклоняется от $m$, все остальные игроки начинают играть $\mu_{-i}^i$ бесконечно долго.
Максимальный однопериодный выигрыш, который может получить игрок $i$ при отклонении:
\begin{equation}
    u_{i,\max} = \max_{a \in A} u_i(a)
\end{equation}
Игрок не захочет отклоняться, если ожидаемая полезность от следования стратегии $\tilde{u}_i = \min_{\tau \ge 0} \bar{u}_i(\tau)$ превышает полезность отклонения:
\begin{equation}
    \tilde{u}_i \ge (1 - \delta) u_{i,\max} + \delta \underline{u}_i \implies \delta \ge \frac{u_{i,\max} - \tilde{u}_i}{u_{i,\max} - \underline{u}_i}
\end{equation}
Поскольку $\tilde{u}_i > \underline{u}_i$ (так как $u \in Y$), всегда найдется $\delta < 1$, удовлетворяющее этому условию.
\end{proof}

Однако данное равновесие опирается на «вечное» наказание, что может быть невыгодно самим наказывающим. Чтобы гарантировать выполнение угроз, требуется концепция совершенства по подыграм.

\begin{theorem}[Народная теорема для СПРН]
Пусть $G = \langle I, A, u \rangle$ — игра в нормальной форме. Пусть $u \in X \cap Y$, и размерность множества $X \cap Y$ равна числу игроков $N$. Тогда существует такое $\bar{\delta} < 1$, что для всех $\delta \in [\bar{\delta}, 1)$ в повторяющейся игре $G(\delta)$ существует совершенное по подыграм равновесие, в котором усредненные дисконтированные выигрыши игроков равны $u$.
\end{theorem}

В отличие от простого равновесия Нэша, стратегия СПРН подразумевает наказание конечной длительности, после которого все игроки возвращаются к профилю действий, дающему каждому строго больше минимаксного выигрыша. Любой игрок, отказывающийся приводить наказание в исполнение, сам подвергается аналогичному конечному наказанию.

\subsection{Пример: Кредитно-денежная политика}

Рассмотрим классическую макроэкономическую модель взаимодействия Центрального банка (ЦБ) и общества (публики).
Экономика описывается уравнением Филлипса:
\begin{equation}
    y = \bar{y} + b(\pi - \pi^e),
\end{equation}
где $y$ — ВВП, $\bar{y}$ — естественный уровень ВВП, $\pi$ — фактическая инфляция, $\pi^e$ — ожидаемая инфляция, $b > 0$. Рост ВВП возможен только при инфляционных сюрпризах ($\pi > \pi^e$).

Функция потерь (отрицательной полезности) публики направлена на точное угадывание инфляции:
\begin{equation}
    u_p = -(\pi - \pi^e)^2
\end{equation}
Функция полезности ЦБ учитывает штраф за инфляцию и стремление достичь целевого ВВП $y^o > \bar{y}$:
\begin{equation}
    u_b = -(y - y^o)^2 - a\pi^2
\end{equation}

В однократной игре (статике) публика первой формирует ожидания $\pi^e$, затем ЦБ выбирает $\pi$. Максимизируя $u_b$ при заданном $\pi^e$, ЦБ выбирает стратегию:
\begin{equation}
    \pi(\pi^e) = \frac{b}{b^2 + a}(y^o - \bar{y} + b\pi^e)
\end{equation}
Публика, понимая стимулы ЦБ, рационально максимизирует $u_p$, что приводит к $\pi^e = \pi$. Подставляя это в реакцию ЦБ, получаем равновесный исход:
\begin{equation}
    \pi^e = \pi = \frac{b}{a}(y^o - \bar{y}), \quad y = \bar{y}
\end{equation}
В результате экономика получает высокую инфляцию без какого-либо реального роста ВВП. Это классическая проблема динамической несостоятельности.

\begin{example}
Построим динамическую модель $G(\delta)$ для $t = 0, 1, 2, \dots$. ЦБ может заявить политику нулевой инфляции $\pi_L = 0$. Альтернатива — обмануть и выбрать инфляционный сюрприз $\pi_H = \frac{b}{b^2 + a}(y^o - \bar{y})$.
Если публика доверяет ЦБ ($\pi^e = 0$), дисконтированный выигрыш ЦБ от честного поведения (инфляция 0 в каждом периоде) составит:
\begin{equation}
    W_{b1} = -\frac{1}{1 - \delta}(y^o - \bar{y})^2
\end{equation}
Если ЦБ отклонится в период $t=0$, он получит временный рост ВВП. Однако, согласно триггерной стратегии публики, начиная с $t=1$ доверие теряется навсегда, и ожидается высокая инфляция $\pi^e = \frac{b}{a}(y^o - \bar{y})$. Выигрыш ЦБ при обмане:
\begin{equation}
    W_{b2} = -(y^o - \bar{y})^2 \frac{a}{b^2 + a} - \frac{\delta}{1 - \delta}(y^o - \bar{y})^2 \left( 1 + \frac{b^2}{a} \right)
\end{equation}
Условие поддержания равновесия с нулевой инфляцией (доверия публики) $W_{b1} \ge W_{b2}$ сводится к ограничению на терпеливость ЦБ:
\begin{equation}
    \delta \ge \frac{a}{2a + b^2}
\end{equation}
Таким образом, низкая инфляция возможна только в повторяющейся игре при достаточно высоком факторе дисконтирования.
\end{example}
